% =====================================================================
\chapter{Atención por producto escalar escalado}
\label{tema:2}
% =====================================================================

Este tema escribe la autoatención de todas las posiciones a la vez como dos productos de matrices y una softmax por filas, resuelve a mano un ejemplo completo, demuestra que la operación es equivariante por permutaciones y justifica el factor $1/\sqrt{d_k}$ con un argumento de varianza y con la jacobiana de la softmax. Se siguen la §3.2.1 de~\cite{vaswani}, las §7.1.1 y §7.3 de~\cite{jm3}, la §11.3 de~\cite{d2l} y la §12.1 de~\cite{bishopdl}.

\section{Forma matricial}

Se apilan las consultas, claves y valores de las $n$ posiciones como filas de $Q=X\WQ\in\R^{n\times d_k}$, $K=X\WK\in\R^{n\times d_k}$ y $V=X\WV\in\R^{n\times d_v}$.

\begin{definicion}[Atención por producto escalar escalado]\label{def:sdpa}
Para $Q\in\R^{n\times d_k}$, $K\in\R^{m\times d_k}$ y $V\in\R^{m\times d_v}$,
\[
\Att(Q,K,V)=\softmax\!\Bigl(\frac{QK\T}{\sqrt{d_k}}\Bigr)V,
\]
donde la softmax se aplica a cada fila. La matriz $A=\softmax(QK\T/\sqrt{d_k})\in\R^{n\times m}$ se llama \emph{matriz} (o \emph{mapa}) \emph{de atención}.
\end{definicion}

La entrada $(i,j)$ de $S=QK\T$ es $q_i\cdot k_j$, de modo que la fila $i$ de $A$ contiene los pesos de la \cref{def:agrupacion} para la consulta $q_i$ con puntuación $q\cdot k/\sqrt{d_k}$, y la fila $i$ de $AV$ es la salida $o_i$. La \cref{tab:formas} resume las formas. Vaswani \emph{et al.} introdujeron el nombre y la fórmula~\cite[§3.2.1, ec.~(1)]{vaswani}; la tarea~1 de CS336 la adopta con $m\ne n$ para cubrir la atención cruzada~\cite[§3.4.4]{cs336}.

\begin{table}[h]
\centering
\begin{tabular}{@{}clll@{}}
\toprule
Paso & Operación & Forma & Significado \\
\midrule
1 & $S=QK\T$ & $n\times m$ & $s_{ij}=q_i\cdot k_j$ \\
2 & $S/\sqrt{d_k}$ & $n\times m$ & escalado (\cref{sec:escalado}) \\
3 & $A=\softmax(S/\sqrt{d_k})$ & $n\times m$ & cada fila es una distribución \\
4 & $O=AV$ & $n\times d_v$ & fila $i$: media ponderada de valores \\
\bottomrule
\end{tabular}
\caption{Formas de las matrices intermedias de la atención.}
\label{tab:formas}
\end{table}

\begin{ejemplo}[Tres tokens]\label{ej:tres}
Sean $d_k=4$, $d_v=2$ y
\[
Q=\begin{pmatrix}2&0&2&0\\0&2&0&2\\2&2&0&0\end{pmatrix},\quad
K=\begin{pmatrix}2&0&0&0\\0&2&0&2\\1&1&1&1\end{pmatrix},\quad
V=\begin{pmatrix}1&0\\0&1\\1&1\end{pmatrix}.
\]
Los parecidos son, por ejemplo, $s_{11}=2\cdot 2=4$ y $s_{22}=2\cdot 2+2\cdot 2=8$; en total, y tras dividir por $\sqrt4=2$,
\[
S=\begin{pmatrix}4&0&4\\0&8&4\\4&4&4\end{pmatrix},\qquad
\frac{S}{2}=\begin{pmatrix}2&0&2\\0&4&2\\2&2&2\end{pmatrix}.
\]
Fila~1: $e^{2}=7{,}389$, $e^{0}=1$, $e^{2}=7{,}389$, de suma $15{,}78$, y pesos $0{,}468$, $0{,}063$, $0{,}468$. Fila~2: $1$, $e^{4}=54{,}60$, $7{,}389$, de suma $62{,}99$, y pesos $0{,}016$, $0{,}867$, $0{,}117$. Fila~3: tres puntuaciones iguales, pesos $1/3$. Por tanto
\[
A=\begin{pmatrix}0{,}468&0{,}063&0{,}468\\0{,}016&0{,}867&0{,}117\\0{,}333&0{,}333&0{,}333\end{pmatrix},\qquad
O=AV=\begin{pmatrix}0{,}937&0{,}532\\0{,}133&0{,}984\\0{,}667&0{,}667\end{pmatrix}.
\]
La consulta~2 busca las coordenadas segunda y cuarta, que ofrece sobre todo la clave~2: la posición~2 se atiende casi solo a sí misma. La consulta~3 es igual de parecida a todas las claves y promedia. Cada fila de $O$ está en la envolvente convexa de las filas de $V$ (\cref{prop:convexa}).
\end{ejemplo}

\begin{figure}[h]
\centering
\begin{tikzpicture}[font=\small,x=11mm,y=9mm]
  \foreach \i/\row in {1/{0.468/0{,}468,0.063/0{,}063,0.468/0{,}468},2/{0.016/0{,}016,0.867/0{,}867,0.117/0{,}117},3/{0.333/0{,}333,0.333/0{,}333,0.333/0{,}333}}{
    \foreach \a/\t [count=\j] in \row {
      \pgfmathsetmacro{\c}{80*\a}
      \fill[black!\c] (\j-1,-\i+1) rectangle (\j,-\i);
      \draw[regla] (\j-1,-\i+1) rectangle (\j,-\i);
      \ifdim\a pt>0.6pt \node[white] at (\j-0.5,-\i+0.5) {$\t$};\else\node at (\j-0.5,-\i+0.5) {$\t$};\fi
    }
    \node[left] at (0,-\i+0.5) {$i=\i$};
  }
  \foreach \j in {1,2,3} \node[above] at (\j-0.5,0) {$j=\j$};
\end{tikzpicture}
\caption{Mapa de atención del \cref{ej:tres}; la intensidad del gris es proporcional a $\alpha_{ij}$. La fila $i$ indica a qué posiciones atiende la posición $i$.}
\label{fig:mapa}
\end{figure}

\section{Equivariancia por permutaciones}

\begin{proposicion}[Equivariancia]\label{prop:equivariancia}
Sea $f(X)=\softmax\bigl(X\WQ(X\WK)\T/\sqrt{d_k}\bigr)X\WV$ la autoatención sin máscara y sea $P\in\R^{n\times n}$ una matriz de permutación. Entonces $f(PX)=P\,f(X)$.
\end{proposicion}

\begin{proof}
Sustituyendo $X$ por $PX$: $Q'=PX\WQ=PQ$, $K'=PK$, $V'=PV$. Luego $S'=PQ(PK)\T=PQK\T P\T=PSP\T$, que reordena filas y columnas de $S$ con la misma permutación. La softmax por filas actúa sobre cada fila por separado y permutar las entradas de un vector permuta del mismo modo las de su softmax; por tanto $A'=PAP\T$. Finalmente $O'=A'V'=PAP\T PV=PAV=PO$, porque $P\T P=I$: las columnas de $P$ son vectores distintos de la base canónica y, por tanto, ortonormales.
\end{proof}

La autoatención no distingue, por sí sola, el orden de los tokens: «el perro muerde al hombre» y «el hombre muerde al perro» producen las mismas salidas, solo que en otras filas. Por eso los modelos de lenguaje añaden información de posición a los vectores de entrada (módulo~14), y por eso la máscara causal del \cref{tema:3}, que sí depende del orden, rompe esta simetría.

\section{Por qué se divide por \texorpdfstring{$\sqrt{d_k}$}{raíz de dk}}
\label{sec:escalado}

Vaswani \emph{et al.} justifican el factor con un argumento de varianza que enuncian en una nota al pie~\cite[§3.2.1, nota~4]{vaswani}; Zhang \emph{et al.}~\cite[§11.3.3]{d2l} y Jurafsky y Martin~\cite[§7.1.1]{jm3} lo reproducen. Se da aquí con demostración.

\begin{teorema}[Varianza del producto escalar]\label{teo:varianza}
Sean $q=(q_1,\dots,q_{d_k})$ y $k=(k_1,\dots,k_{d_k})$ vectores aleatorios cuyas $2d_k$ componentes son independientes, de media~0 y varianza~1. Entonces
\[
\E[q\cdot k]=0,\qquad \Var(q\cdot k)=d_k,\qquad \Var\Bigl(\frac{q\cdot k}{\sqrt{d_k}}\Bigr)=1 .
\]
\end{teorema}

\begin{proof}
Sea $X_i=q_ik_i$. Por la independencia de $q_i$ y $k_i$, $\E[X_i]=\E[q_i]\,\E[k_i]=0$. Además $\E[q_i^{2}]=\Var(q_i)+\E[q_i]^{2}=1$, e igual para $k_i$; como $q_i^{2}$ y $k_i^{2}$ también son independientes, $\Var(X_i)=\E[q_i^{2}k_i^{2}]-\E[X_i]^{2}=\E[q_i^{2}]\,\E[k_i^{2}]=1$. Las variables $X_1,\dots,X_{d_k}$ son independientes entre sí, porque cada una es función de un par de componentes distinto, de modo que la varianza de su suma es la suma de las varianzas (módulo~5): $\Var(q\cdot k)=\sum_i\Var(X_i)=d_k$. Por linealidad, $\E[q\cdot k]=0$, y $\Var(cY)=c^{2}\Var(Y)$ con $c=1/\sqrt{d_k}$ da la última igualdad.
\end{proof}

Por el teorema central del límite, $q\cdot k/\sqrt{d_k}$ es aproximadamente $\mathcal N(0,1)$ para $d_k$ grande. Sin el factor, con $d_k=64$ los parecidos tendrían desviación típica~8 en la inicialización y diferencias de~16 entre dos claves serían frecuentes. Una simulación con $2\cdot 10^{5}$ pares de vectores gaussianos da varianzas empíricas $4{,}00$ para $d_k=4$ y $63{,}8$ para $d_k=64$.

Para ver por qué importa la escala se necesita la jacobiana de la softmax, calculada en el módulo~4 con la regla del cociente.

\begin{proposicion}[Saturación de la softmax]\label{prop:saturacion}
Sea $p=\softmax(z)\in\R^{n}$ y $J=\partial p/\partial z$, de entradas $J_{ij}=p_i(\delta_{ij}-p_j)$. Entonces
\[
\max_{i,j}\abs{J_{ij}}\;\le\;1-\max_{l}p_l .
\]
En consecuencia, si una probabilidad tiende a~1, todas las derivadas parciales de la softmax tienden a~0.
\end{proposicion}

\begin{proof}
Sea $p_k=\max_lp_l$, de modo que $p_i\le 1-p_k$ para todo $i\ne k$ (las probabilidades suman~1). Si $i=j$, $\abs{J_{ii}}=p_i(1-p_i)$: para $i=k$ vale $p_k(1-p_k)\le 1-p_k$, y para $i\ne k$, $p_i(1-p_i)\le p_i\le 1-p_k$. Si $i\ne j$, $\abs{J_{ij}}=p_ip_j$ y al menos uno de los dos índices es distinto de $k$; si lo es $i$, $p_ip_j\le p_i\le 1-p_k$, y análogamente si lo es $j$.
\end{proof}

\begin{table}[h]
\centering
\begin{tabular}{@{}lll@{}}
\toprule
Logits $z$ & $\softmax(z)$ & $\max\abs{J_{ij}}$ \\
\midrule
$(1,0,0)$ & $(0{,}576;\allowbreak\ 0{,}212;\allowbreak\ 0{,}212)$ & $0{,}244$ \\
$(8,0,0)$ & $(0{,}9993;\allowbreak\ 0{,}0003;\allowbreak\ 0{,}0003)$ & $0{,}00067$ \\
\bottomrule
\end{tabular}
\caption{Con logits ocho veces mayores, las derivadas de la softmax se reducen en un factor de unos 360.}
\label{tab:saturacion}
\end{table}

\begin{observacion}[El factor como temperatura]
Dividir por $\sqrt{d_k}$ equivale a aplicar la softmax con temperatura $\tau=\sqrt{d_k}$ (módulo~1, tema~4). Sin ella, el modelo empezaría en la región saturada: en el \cref{ej:tres}, sin dividir por~2, el peso de la fila~2 sobre la posición~2 pasa de $0{,}867$ a $0{,}982$. Vaswani \emph{et al.} observan, en la misma línea, que sin escalado la atención aditiva supera a la de producto escalar cuando $d_k$ es grande~\cite[§3.2.1]{vaswani}. Durante el entrenamiento, la red puede modificar la escala efectiva agrandando o encogiendo $\WQ$ y $\WK$.
\end{observacion}

\begin{observacion}[Evaluación estable]
Como la softmax es invariante al restar una constante a todas las entradas, se evalúa restando el máximo de cada fila antes de exponenciar, lo que evita desbordamientos en fp16 (módulo~11). La tarea~1 de CS336 exige implementarla así~\cite[§3.4.4]{cs336}.
\end{observacion}

\begin{lecturas}
\item Vaswani \emph{et al.}~\cite{vaswani}, §3.2.1 «Scaled Dot-Product Attention» y su nota~4.
\item Jurafsky y Martin~\cite{jm3}, §7.1.1 «Attention more formally» y §7.3 «Parallelizing computation using a single matrix X».
\item Zhang \emph{et al.}~\cite{d2l}, §11.3 «Attention Scoring Functions», en particular §11.3.3.
\item Bishop y Bishop~\cite{bishopdl}, cap.~12, §12.1 «Attention».
\end{lecturas}

\section*{Ejercicios del tema 2}
\addcontentsline{toc}{section}{Ejercicios del tema 2}

\begin{ejercicio}{1}{eje:2-formas}
¿Qué formas tienen $S=QK\T$, $A$ y $O$ en la autoatención si $n=10$, $d_k=16$ y $d_v=32$? ¿Por qué $d_k$ no aparece en la forma de $O$?
\end{ejercicio}

\begin{ejercicio}{2}{eje:2-var}
En las hipótesis del \cref{teo:varianza}, calcule $\Var(q\cdot k)$ y la desviación típica de $q\cdot k$ y de $q\cdot k/\sqrt{d_k}$ para $d_k=64$ y $d_k=128$.
\end{ejercicio}

\begin{ejercicio}{2}{eje:2-temp}
Los logits de cuatro palabras candidatas son $z=(1{,}2;\allowbreak\ 0{,}9;\allowbreak\ 0{,}1;\allowbreak\ -0{,}5)$. Calcule $\softmax(z/\tau)$ para $\tau=0{,}2$ y $\tau=2$, compruebe que cada distribución suma~1 y describa los límites $\tau\to 0^{+}$ y $\tau\to\infty$. Relacione el resultado con el factor $1/\sqrt{d_k}$.
\fuente{adaptado de Jurafsky y Martin~\cite{jm3}, cap.~7, ejercicio~7.5 (logits propios).}
\end{ejercicio}

\begin{ejercicio}{2}{eje:2-jac}
Para $z=(1,0,0)$, calcule la matriz jacobiana completa de la softmax y compruebe la cota de la \cref{prop:saturacion}. Verifique que la suma de cada columna de $J$ es~0 y explique por qué debe serlo.
\end{ejercicio}

\begin{ejercicio}{3}{eje:2-sigma}
Generalice el \cref{teo:varianza} al caso en que las componentes tienen media~0 y varianza $\sigma^{2}$. ¿Qué factor $c$ hace que $\Var(c\,q\cdot k)=1$?
\end{ejercicio}

\begin{ejercicio}{3}{eje:2-cruzada}
En la atención cruzada, las consultas $Q=Y\WQ$ proceden de una secuencia $Y\in\R^{m\times d}$ y las claves y valores, $K=X\WK$ y $V=X\WV$, de otra $X\in\R^{n\times d}$. Demuestre que la salida es \emph{invariante} (no solo equivariante) si se permutan las filas de $X$: para toda permutación $P$, $\Att(Q,PK,PV)=\Att(Q,K,V)$.
\end{ejercicio}
